\PassOptionsToPackage{table}{xcolor}
\documentclass[journal, twocolumn]{IEEEtran}

\usepackage{amsmath,amsfonts}
\usepackage{array}
\usepackage{textcomp}
\usepackage{stfloats}
\usepackage{url}
\usepackage{verbatim}
\usepackage{graphicx}
\usepackage{subfiles}
\usepackage{subcaption}
\usepackage{booktabs}
\usepackage{multirow}
\usepackage{makecell}
\usepackage{algorithm}
\usepackage{algpseudocode}
\usepackage{enumitem}
\usepackage{fmtcount}
\usepackage{tcolorbox}
\usepackage{cite}
\usepackage{amssymb}
\usepackage{xcolor}
\usepackage{hyperref}
\usepackage{tabularx}
\usepackage{amsthm}
\usepackage{pifont}
\usepackage{balance}
\usepackage{tikz}

\newcommand{\cmark}{\ding{51}}%
\newcommand{\xmark}{\ding{55}}%
\newtheorem{theorem}{Theorem}[section]
\newtheorem{proposition}[theorem]{Proposition}
\newtheorem{definition}[theorem]{Definition}
\newtheorem{lemma}[theorem]{Lemma}
\newtheorem{corollary}[theorem]{Corollary}
\newtheorem{property}{Property}
\newtheorem{problem}{Problem}
\newtheorem{invariant}{Invariant}
\newtheorem*{remark}{Remark}

\usetikzlibrary{shapes.geometric, arrows.meta, positioning, fit, backgrounds, calc, decorations.pathreplacing, shapes.multipart}

% Graceful handling of missing figure files: show placeholder box with filename
\let\origincludegraphics\includegraphics
\renewcommand{\includegraphics}[2][]{%
  \IfFileExists{#2}{%
    \origincludegraphics[#1]{#2}%
  }{\IfFileExists{#2.pdf}{%
    \origincludegraphics[#1]{#2}%
  }{\IfFileExists{#2.png}{%
    \origincludegraphics[#1]{#2}%
  }{\IfFileExists{#2.jpg}{%
    \origincludegraphics[#1]{#2}%
  }{\IfFileExists{#2.jpeg}{%
    \origincludegraphics[#1]{#2}%
  }{%
    \fbox{\parbox{0.8\linewidth}{\centering\small\ttfamily [Missing figure: \detokenize{#2}]}}%
  }}}}}%
}

\hyphenation{op-tical net-works semi-conduc-tor IEEE-Xplore rule-book}
\def\BibTeX{{\rm B\kern-.05em{\sc i\kern-.025em b}\kern-.08em
    T\kern-.1667em\lower.7ex\hbox{E}\kern-.125emX}}


\begin{document}

\title{Closed-Loop Deployment of Lexicographic Constraint Programming for Trajectory Model Predictive Control on the nuPlan Simulator}

\author{Hadi~Hajieghrary%
\thanks{H.~Hajieghrary is with Torc Robotics, Blacksburg, VA, USA.
        E-mail: \texttt{Hadi.Hajieghrary@Torc.ai}.}%
\thanks{The author thanks the nuPlan development team for the simulator infrastructure used in this study.}%
}

\maketitle

\begin{abstract}
We are interested in the closed-loop deployment of priority-ordered constrained Model Predictive Control (MPC) for autonomous-vehicle trajectory planning. Building on the geometric equivalence framework of Lexicographic Constraint Programming (LCP)~\cite{lcp2025}, we present (i) a dynamics-agnostic reference implementation of the framework's Algorithms~0, 1A, and 1B together with its pre-flight diagnostics (\S8.5), online cascade-fallback policy (\S9.6), and Relaxation Decision Framework (\S10), verified to $10^{-6}$ precision against the published numerical answers of the framework's worked examples by a 47-test regression suite; (ii) a closed-loop deployment, on the nuPlan simulator, that realises the calibrated weighted-sum substitution of the formal lex cascade through a two-level MPC pipeline with eleven implementation primitives the convex framework does not directly address; and (iii) a dual-condition demonstration protocol on the 16-scenario \texttt{nuPlan-mini} benchmark in which the operational weighted-sum planner ($C_1$) and the formal lex cascade ($C_4$) are run on each scenario. The empirical evidence confirms priority-respecting compromise on every scenario (rule violations concentrate at the low-priority comfort and headway levels while high-priority collision and traffic-light levels remain effectively held), reproduces the formal cascade's per-tick compliance pattern at the high-priority levels (the operational realisation of Theorem~4.1 of~\cite{lcp2025}), and quantifies the architectural wall-time benefit of substituting the calibrated weighted-sum solve for the $L{+}1$-stage cascade. The framework's standing assumptions (FM~I LICQ + FM~II convexity, \cite{lcp2025} \S8.5) hold across the benchmark, and the empirical \S10.2 necessity scan identifies the scenarios on which the deployment legitimately operates in the necessity-required relaxation regime --- ascribing the highest-level violations to road geometry rather than to planner suboptimality. The deployment and the protocol are released as open artefacts.
\end{abstract}

\begin{IEEEkeywords}
Lexicographic optimization, model predictive control, weighted-sum scalarization, rulebook, autonomous vehicles, convex optimization, priority-ordered constraints, closed-loop simulation, nuPlan, compliance vector, support-normal equivalence.
\end{IEEEkeywords}


\section{Introduction}
\label{sec:introduction}
\subfile{./Sections/Section_I_Introduction}

\section{Foundations: The LCP Framework}
\label{sec:foundations}
\subfile{./Sections/Section_II_Foundations}

\section{The 25-Rule Rulebook and Its Partition}
\label{sec:rulebook}
\subfile{./Sections/Section_III_Rulebook}

\section{Closed-Loop Deployment: System Architecture}
\label{sec:deployment}
\subfile{./Sections/Section_IV_Deployment}

\section{Demonstration Protocol}
\label{sec:protocol}
\subfile{./Sections/Section_V_Protocol}

\section{Single-Condition System Behaviour}
\label{sec:singlecondition}
\subfile{./Sections/Section_VI_SingleConditionBehavior}

\section{Empirical Results}
\label{sec:results}
\subfile{./Sections/Section_VII_Results}

\section{Discussion}
\label{sec:discussion}
\subfile{./Sections/Section_VIII_Discussion}

\section{Conclusion and Future Work}
\label{sec:conclusion}
\subfile{./Sections/Section_IX_Conclusion}


\bibliographystyle{./IEEEtran}
\bibliography{./References/references.bib}


\appendices

\section{Hyperparameter Reference}
\label{app:hyperparameters}
\subfile{./Sections/Appendix_A_Hyperparameters}

\section{Algorithm Pseudocode (0, 1A, 1B)}
\label{app:algorithms}
\subfile{./Sections/Appendix_B_AlgorithmPseudocode}

\section{Encoder Specification}
\label{app:encoders}
\subfile{./Sections/Appendix_C_EncoderSpec}

\section{Framework Verification}
\label{app:verification}
\subfile{./Sections/Appendix_D_FrameworkVerification}

\section{Reproducibility Manifest}
\label{app:reproducibility}
\subfile{./Sections/Appendix_E_Reproducibility}

\balance

\end{document}
